% Data. States no result. Sources: plan.json design.data and design.method (OUTCOMES, SEASONAL
% ADJUSTMENT), data/CODEBOOK.md (counts and checks only), sources/notes.md sections 5, 9-12.
% Summary statistics: Table tab:sumstats from the r-desc-* claims of run describe (deviation D8,
% added after paper-r1 P3; levels only, unreviewed with the human's agreement).
\section{Data}
\label{sec:data}

\subsection{Employment and wages: the QCEW}

Employment and wages come from the Quarterly Census of Employment and Wages (QCEW) of the U.S.
Bureau of Labor Statistics, through its open-data files by industry
\citep{blsqcewopendata}. We use the quarters 2014Q1--2026Q1 for four industries: private NAICS
722513 (limited-service restaurants), private 722511 (full-service restaurants), private 7225
(restaurants and other eating places) and industry 10 (total covered employment, all ownerships):
196 files of one industry and quarter each. On November 21, 2026 the files for 2026Q2 were not yet
available; the BLS schedules them for December 2, 2026 \citep{bls2026cewq1}, and we do not wait
for them. The vintage used is pinned by the SHA-256 hash of every file in the data manifest. The
files for 2025Q1--2026Q1 carry a server date of August 21, 2026. Under the BLS revision schedule
a year's quarters receive their last revision with the release of the following year's
first-quarter data, so in this vintage, the 2026Q1 release, all four 2025 quarters are final and
only 2026Q1 is preliminary; 2026Q1 will be revised four more times, first with the 2026Q2 release
\citep{blsqcewrevisions,bls2026cewq1}. A later download would be a new vintage. The data are in
the public domain.

We use state rows only: aggregation level 58 for 722513 and 722511, 56 for 7225 and 50 for
industry 10, with ownership code 5 (private) for the restaurant industries and 0 (total covered)
for industry 10 \citep{blsqcewagglvl,blsqcewowncodes}. The units are the 50 states and the
District of Columbia; Puerto Rico and the Virgin Islands are excluded. QCEW monthly employment
counts the workers who worked or received pay in the pay period including the 12th of the month,
counting a person with two jobs once for each employer, and the average weekly wage is total
quarterly wages divided by mean monthly employment and by 13
\citep{blsqcewhom,blsqcewtechnote}. Employment is therefore a count of jobs, not of hours, and the
weekly wage is earnings per job, not the hourly wage the law sets. Quarterly employment is the
mean of the three monthly levels. Wages are nominal.

No state row of 722513 or 722511 is suppressed in any quarter, and state totals include the data
of suppressed county cells \citep{blsqcewfaq}. Four state rows of 7225 are suppressed, all in
2024Q4: Alaska, New Hampshire, Vermont and Wyoming. We do not use counties. In the county files
the share of suppressed restaurant cells rises in 2025Q1 in all three restaurant industries and
stays higher through 2026Q1; for 7225 it goes from 23.2 percent of county cells in 2024Q4 to 29.6
percent in 2025Q1, and from 2024 to 2025 the number of suppressed 7225 county cells rises in 45
of the 53 state-level areas. A county panel would therefore be selected partly on disclosure in
the post-period. We never use NAICS 722515, where covered chains are also coded
(Section~\ref{sec:coverage}).

\subsection{Population}

Resident population on July 1 comes from the Census Bureau's Population Estimates Program: the
2010--2020 intercensal state estimates for 2013--2019, summed over race and age groups, and the
Vintage 2025 estimates for 2020--2025 \citep{censuspep2020int,censuspepv2025}. The log of
population is interpolated linearly between July 1 anchors, placed at mid-year, with quarter $q$
of year $y$ placed at $y + (q-0.5)/4$; after mid-2025 it is extrapolated with the 2024--2025 log
growth, so 2025Q3, 2025Q4 and 2026Q1 are extrapolated and flagged. At the splice of the two files
no state's ratio of the July 2020 Vintage 2025 estimate to the July 2019 intercensal estimate
lies outside 0.97--1.03. The 2019Q4 value is interpolated between the July 2019 and July 2020
anchors. Population enters the per-capita outcome, the population weights, the scaling of one of
our two placebo standard errors and the ranking of states by size (Section~\ref{sec:inference});
the extrapolated quarters matter only if a per-capita candidate is selected.

\subsection{Variables and outcomes}

Table~\ref{tab:variables} lists the variables of the state panel. From them the estimator code
builds three outcomes for California and for each donor or donor aggregate:
\begin{itemize}
  \item $E$, the log of quarterly 722513 employment;
  \item $P$, seasonally adjusted log employment minus log population (employment per capita);
  \item $W$, the log of the average weekly wage as the BLS publishes it, and for a donor
        aggregate the log of summed quarterly wages divided by 13 times summed mean employment.
\end{itemize}
Every series that enters an estimator is seasonally adjusted after aggregation, with fixed
factors estimated on 2014--2019. A centred $2\times4$ moving average
$m_t = (0.5y_{t-2} + y_{t-1} + y_t + y_{t+1} + 0.5y_{t+2})/4$ is computed for 2014Q3--2019Q2, so
that it uses only 2014Q1--2019Q4; the factor $s_q$ of quarter type $q$ is the mean of $y_t - m_t$
over its five quarters, minus the mean of the four factors; and the adjusted series is
$y_t - s_{q(t)}$ for every $t$. The same factors serve every run, including the runs that read
only data through 2019Q4.

\begin{table}[tbp]
\centering
\caption{Variables of the state panel}
\label{tab:variables}
\small
\begin{tabularx}{\linewidth}{@{}l>{\raggedright\arraybackslash}X@{}}
\toprule
Variable & Definition \\
\midrule
\texttt{emp\_722513} & Mean of the three monthly employment levels, private 722513, state row \\
\texttt{emp\_722511}, \texttt{emp\_7225} & The same for private 722511 and private 7225 (missing on the four suppressed 7225 rows) \\
\texttt{emp\_10} & The same for industry 10, total covered employment, private and government \\
\texttt{emp\_nonff} & \texttt{emp\_10} minus \texttt{emp\_722513}: employment outside limited-service restaurants \\
\texttt{sep\_722513} & Third-month employment of 722513, for the cross-check against the fixed-horizon estimates of CEMN \\
\texttt{wages\_722513} & Total quarterly wages of 722513, nominal dollars \\
\texttt{aww\_722513} & Average weekly wage of 722513 as published; equal to wages divided by 13 times employment within \$1 in every row \\
\texttt{pop} & July 1 resident population interpolated to the quarter; \texttt{pop\_extrap} flags 2025Q3--2026Q1 \\
\bottomrule
\end{tabularx}
\end{table}

The full panel has 2,499 rows, 51 units by 49 quarters, with no missing value of 722513, 722511 or
industry 10. A second file holds the same rows for 2014Q1--2019Q4 only (1,224 rows); it is the
only input of the two runs that select the primary estimator and check the level of the inference
(Sections~\ref{sec:select} and~\ref{sec:inference}). Both files come from one data-preparation run
that computes no statistic comparing states or periods.

\subsection{Summary statistics}
\label{sec:sumstats}

Table~\ref{tab:sumstats} gives the levels of private 722513 employment, average weekly wage and
establishments for California and the two donor pools, over the fit window and the post window.
It is descriptive only: the run that produced it, \cand{describe}, was added at the paper stage
after the paper review asked for the table (deviation D8), computes no statistic comparing
California with another unit or one period with another, and was run without a code review with
the human PI's agreement, because the code review rounds had been used. California's mean
employment is several times that of the average donor state in either pool, and its weekly wage
is higher in both windows; this is why the estimators work with logs and why California's level
of log employment lies outside the donors' convex hull (Section~\ref{sec:candidates}). The table
reports means, not dispersion or population, which the run does not compute. The two California
levels used for conversions elsewhere are the post-window mean employment in the table and the
2023Q3 average weekly wage, \$$\AwwCaBase$.

\begin{table}[tbp]
\centering
\caption{Summary statistics: private NAICS 722513, California and the donor pools}
\label{tab:sumstats}
\footnotesize
\setlength{\tabcolsep}{4pt}
\begin{tabular}{@{}lrrrrrr@{}}
\toprule
& \multicolumn{2}{c}{California} & \multicolumn{2}{c}{R50} & \multicolumn{2}{c}{D20} \\
\cmidrule(lr){2-3}\cmidrule(lr){4-5}\cmidrule(l){6-7}
& Fit & Post & Fit & Post & Fit & Post \\
\midrule
Employment per state, mean & \DescCAEmpMeanPre & \DescCAEmpMeanPost & \DescRFiftyEmpMeanPre & \DescRFiftyEmpMeanPost & \DescDTwentyEmpMeanPre & \DescDTwentyEmpMeanPost \\
Employment, pool total & & & \DescRFiftyEmpTotalPre & \DescRFiftyEmpTotalPost & \DescDTwentyEmpTotalPre & \DescDTwentyEmpTotalPost \\
Average weekly wage (\$) & \DescCAAwwPre & \DescCAAwwPost & \DescRFiftyAwwPre & \DescRFiftyAwwPost & \DescDTwentyAwwPre & \DescDTwentyAwwPost \\
Establishments per state, mean & \DescCAEstMeanPre & \DescCAEstMeanPost & \DescRFiftyEstMeanPre & \DescRFiftyEstMeanPost & \DescDTwentyEstMeanPre & \DescDTwentyEstMeanPost \\
Establishments, pool total & & & \DescRFiftyEstTotalPre & \DescRFiftyEstTotalPost & \DescDTwentyEstTotalPre & \DescDTwentyEstTotalPost \\
\addlinespace
States & \DescCANStatesPre & \DescCANStatesPost & \DescRFiftyNStatesPre & \DescRFiftyNStatesPost & \DescDTwentyNStatesPre & \DescDTwentyNStatesPost \\
Quarters & \DescCANQuartersPre & \DescCANQuartersPost & \DescRFiftyNQuartersPre & \DescRFiftyNQuartersPost & \DescDTwentyNQuartersPre & \DescDTwentyNQuartersPost \\
State-quarters & \DescCANObsPre & \DescCANObsPost & \DescRFiftyNObsPre & \DescRFiftyNObsPost & \DescDTwentyNObsPre & \DescDTwentyNObsPost \\
\bottomrule
\end{tabular}

\medskip
\begin{minipage}{\linewidth}\footnotesize
\emph{Notes:} Run \cand{describe} (deviation D8; descriptive, levels only; run without a code
review with the human PI's agreement). Fit window: 2014Q1--2019Q4 and 2022Q1--2023Q3 (31
quarters); post window: 2024Q2--2026Q1 (8 quarters; 2026Q1 preliminary). Employment and
establishments: quarterly means of the published, not seasonally adjusted state values, per state
or summed over the pool. Average weekly wage: summed quarterly wages divided by 13 times summed
employment, nominal dollars. R50: the 50 other units; D20: the 20 federal-minimum states.
California is in neither pool.
\end{minipage}
\end{table}

\subsection{Comparison pools}
\label{sec:pools}

California (FIPS 06) is the treated unit. Each candidate estimator uses one of two donor pools.
D20 is Dube's set of 20 states whose binding fast-food minimum is the federal \$7.25: Alabama,
Georgia, Idaho, Indiana, Iowa, Kansas, Kentucky, Louisiana, Mississippi, New Hampshire, North
Carolina, North Dakota, Oklahoma, Pennsylvania, South Carolina, Tennessee, Texas, Utah, Wisconsin
and Wyoming \citep{dube2026labor}. In each of them the binding minimum for employers covered by
the Fair Labor Standards Act was \$7.25 throughout 2014--2019 \citep{dolwhd2025history,dolwhd2026state}.
R50 is all 50 other units, the ``rest of the US'' of Clemens and coauthors, here including Alaska
and Michigan, which CEMN drop from their monthly panel; our quarterly state rows have no gaps.
Most R50 units outside D20 raised their own minimum wage during 2014--2019
(Section~\ref{sec:ramp}).

\subsection{Other data}

Two further sources enter no estimate of the study. The Proposition 99 panel distributed with the
\texttt{synthdid} R package, a long-format version of the data of \citet{abadie2010synthetic} for
39 states and 1970--2000 \citep{synthdidpackage}, is used only to check our estimator code against
the published values of \citet[Table~1]{arkhangelsky2021synthetic}. The Wage and Hour Division's
history of state minimum wages \citep{dolwhd2025history} is used only to confirm that every D20
state was at the federal minimum in 2014--2019, which defines the set of states that calibrates our
critical values (Section~\ref{sec:inference}).

\subsection{Software}

The study runs in Python 3.14.6 with numpy 2.5.3, scipy 1.18.1, pandas 3.0.5 and statsmodels
0.15.0, with the environment frozen in the replication package. Nothing in a study run is random.
Weights are solved by sequential least squares on the simplex, starting from uniform weights, and
their Karush--Kuhn--Tucker conditions are checked to $10^{-6}$.
